Quotes in \hyperref[meth:carry-trade-uip]{the carry trade and UIP} follow the convention in \hyperref[def:fx-market-real-exchange-rates]{the FX market}, and its numbers are checked in \nolinkurl{trading_notes.numerical.fx_mechanics_parity.fx_mechanics_checks}.

\begin{qfmethod}[The carry trade and uncovered interest parity: the trade's construction and profit, why it profits when UIP fails, crash risk, negative skew and leverage, and safe-haven funding currencies]\label{meth:carry-trade-uip}
\qffield{Statement} Under UIP high-yield currencies should depreciate by the interest differential, so borrowing low-yield and investing high-yield currencies would earn nothing on average. ``Overwhelmingly, the evidence suggests that uncovered interest rate parity does not hold, at least over short- and medium-run time periods'': high-yield currencies have not depreciated, and low-yield currencies have not appreciated, as predicted \citep[chap.~10, sec.~4, pp.~569--572]{economics-investment-decision-makers-2013}. The \emph{carry trade} is long high-yield and short low-yield (\emph{funding}) currencies. Its reward is ``the gradual accrual of the interest rate differential''; its risk is an adverse exchange-rate move that ``can occur virtually instantaneously'' (p.~570). Returns are more peaked than normal, with fat tails and negative skew: many small gains and occasional large losses. Crashes come because the trade is leveraged and crowded in calm periods, unwinds trigger stop-losses and further liquidation, and in turmoil investors flee to low-yield havens such as the yen and Swiss franc (pp.~570--571). One explanation is that the yield is a risk premium for less stable economies; Menkhoff, Sarno, Schmeling and Schrimpf find that ``high interest rate currencies are negatively related to innovations in global FX volatility and thus deliver low returns in times of unexpected high volatility, when low interest rate currencies provide a hedge'', and that their volatility proxy captures more than 90 percent of the cross-sectional excess returns of five carry portfolios \citep[abstract]{msss-2012}.

\qffield{Notation} Funded P\&L per unit of funding: $\Pi=(S_T/S_0)(1+i_f\tau_f)-(1+i_d\tau_d)$, with $S$ in funding currency per unit of target (\hyperref[mod:fx-carry-trade]{FX carry}).

\qffield{Conditions} The profit is an expected return in normal periods, not an arbitrage; it is earned by bearing crash risk. Transaction costs and the forward points actually traded (including any basis) reduce it.

\qffield{Practice} \emph{Original.}
\begin{itemize}
\item Borrowing yen at 0.1 percent and investing in Australian dollars at 4.5 percent for a year earns 1.265 percent of the funding if the Australian dollar falls 3 percent: the 4.4-point carry less the depreciation and a 0.135-point cross term. UIP would have predicted a 4.21 percent fall, which would have left nothing.
\item A simulated monthly stream earning 0.6 percent in normal months but losing 8 percent in 2 percent of months has a positive mean (0.45 percent) and a skewness of $-1.9$.
\end{itemize}

\qffield{Pitfalls} Judging the trade by its Sharpe ratio in calm years. Ignoring that funding currencies rally in crises, when other risky assets fall. Assuming diversification across target currencies removes the crash risk, which is common to them.

\qffield{See also} \hyperref[mod:fx-carry-trade]{FX carry model}; \hyperref[res:international-parity-conditions]{UIP}; \hyperref[meth:active-currency-strategies]{active currency strategies}.

\end{qfmethod}

\begin{qfmodel}[The FX carry trade: carry as the rate differential at a constant exchange rate, funding and target currencies, crowding and crash risk, and carry-to-risk ratios]\label{mod:fx-carry-trade}
\qffield{Statement} ``In currencies, carry is the return one makes on holding a currency, assuming that the exchange rate is held constant''; carry trades borrow in low-yielding currencies (yen, dollar, sterling in Gliner's time) and buy high-yielding ones (Australian and New Zealand dollars, Brazilian real, Turkish lira) \citep[chap.~7, sec.\ ``Carry Indicators for Currency Valuation'', PDF pp.~206--207]{s7-gliner}. ``The carry trade is a very crowded trade, which means that in large risk off markets such as 2008, the strategy suffers big losses'' (p.~207). \emph{Risk-adjusted carry} divides three-month carry by 30-day volatility; Gliner's Table~7.6 (September 2013) ranks pairs within groups, NZD/EUR at 0.26 being the most attractive New Zealand cross, and he says risk-adjusted carry ``tends to outperform standard carry trades'' and needs frequent adjustment (pp.~208--209).

\qffield{Notation} \emph{Conventions}. $S_t$ is the spot rate quoted as domestic currency per unit of foreign currency; $i_d$ and $i_f$ are simple money-market rates, each quoted on its own day-count basis, and $\tau_d$ and $\tau_f$ are the accrual fractions for the same calendar period under those bases (for 90 days, $90/360$ for an ACT/360 rate and $90/365$ for an ACT/365 rate); $F$ is the forward for the period. \emph{Exact relations}: covered interest parity,
\[ \frac{F}{S_0}=\frac{1+i_d\tau_d}{1+i_f\tau_f}; \]
and the funded long-foreign position (borrow one unit of domestic currency at $i_d$, buy $1/S_0$ of foreign currency, deposit it at $i_f$, convert back at $S_T$) earns, per initial domestic unit, the arithmetic P\&L
\[ \Pi=\frac{S_T}{S_0}(1+i_f\tau_f)-(1+i_d\tau_d). \]
Carry is $\Pi$ at an unchanged rate, $S_T=S_0$: exactly $i_f\tau_f-i_d\tau_d$, which equals $(i_f-i_d)\tau$ only when both rates share one basis. Hedging at $F$ gives $\Pi=(F/S_0)(1+i_f\tau_f)-(1+i_d\tau_d)=0$ exactly, on either basis. \emph{First-order forms}: with $s=S_T/S_0-1$, $\Pi=(i_f\tau_f-i_d\tau_d)+s+s\,i_f\tau_f$, so $\Pi\approx i_f\tau_f-i_d\tau_d+s$ after dropping the cross term; the forward premium $(F-S_0)/S_0=(i_d\tau_d-i_f\tau_f)/(1+i_f\tau_f)\approx i_d\tau_d-i_f\tau_f$.

\emph{A different quantity}: the continuously compounded log excess return $x=\ln(S_T/S_0)+\ln(1+i_f\tau_f)-\ln(1+i_d\tau_d)$, equivalently $\ln(S_T/S_0)+(r_f-r_d)\tau$ with continuously compounded rates on a common $\tau$; it satisfies $e^x-1=\Pi/(1+i_d\tau_d)$, a return per unit of domestic funding repaid, not the arithmetic P\&L $\Pi$.

\emph{Example} (90 days, $i_d=5\%$ ACT/360, $i_f=2\%$ ACT/365): $\tau_d=0.25$, $\tau_f=0.246575$, $F/S_0=1.0075314$; carry $-0.75685\%$, not the $-0.75\%$ a common basis would give; a 1 percent rise in $S$ gives $\Pi=+0.24808\%$ against $+0.24315\%$ to first order, and $x=0.24472\%$. The carry-to-risk ratio divides annualised carry by annualised volatility.

\qffield{Conditions} The realised arithmetic P\&L is carry plus the exchange-rate change, to first order (above); uncovered interest parity predicts that the change offsets carry on average, and carry's profitability is the forward-premium puzzle. Carry-to-risk uses volatility as the risk measure and so ignores skewness and crash risk; the ranking by ratio can differ from the ranking by carry.

\qffield{Practice} The usual carry risks (crowding, weekend gaps in EM currencies, policy regime shifts, correlated sell-offs of carry currencies against funders, illiquidity when exits are needed) are consistent with Gliner's 2008 remark but go beyond it. That exchange-rate moves often dominate carry returns, and that carry has positive average returns with significant negative skewness, are findings of the carry literature. The abstract of Brunnermeier, Nagel and Pedersen's working paper states that ``exchange rate movements between high-interest-rate and low-interest-rate currencies are negatively skewed'', attributes the skewness to ``sudden unwinding of carry trades'' when risk appetite and funding liquidity fall, and reports that carry losses reduce future crash risk but raise its price \citep[abstract, NBER Working Paper 14473, November 2008, revised December 2013]{catalogue-0241}.

\emph{The paper's body}: 3-month interbank rates and exchange rates against the dollar for eight developed currencies (AUD, CAD, JPY, NZD, NOK, CHF, GBP, EUR) from Datastream, 1986--2006 (PDF p.~8); equal-weighted portfolios long the $k$ highest-rate and short the $k$ lowest-rate currencies, rebalanced weekly (1992--2006) or quarterly (1986--2006), have skewness of $-0.717$ and $-0.700$ for $k=1$, $-0.537$ and $-0.748$ for $k=2$ and $-0.695$ and $-0.977$ for $k=3$ (weekly and quarterly), with annualised Sharpe ratios between 0.59 and 0.78 and average quarterly returns of 1.6 to 2.2 percent (Table~2, PDF p.~12); the CRSP value-weighted U.S.\ stock index has skewness of $-0.41$ weekly and $-0.88$ quarterly over the same samples (p.~13), so carry skewness is ``roughly comparable'' to the stock market's, and it is not diversified away by adding currencies \citep[working paper of November 2008, PDF pp.~8 and 12--13]{catalogue-0241}.

These are developed-market G10 samples to 2006, before the 2008 unwind. \emph{A broader sample through the crisis} (authors' accepted version, February 2011). Menkhoff, Sarno, Schmeling and Schrimpf sort up to 48 currencies against the dollar (and a subsample of 15 developed ones, ten after the euro) into five portfolios by one-month forward discount each month, November 1983 to August 2009, net of bid--ask costs charged on entry and exit with about 30 percent monthly turnover \citep[sec.~II, pp.~8--9]{msss-2012}.

The high-minus-low carry portfolio earned 7.23 percent a year net with a standard deviation of 9.81 percent, skewness $-1.03$, kurtosis 4.79 and a Sharpe ratio of 0.74 for all countries, and 5.72 percent, 10.24, $-0.92$ and 0.56 for developed countries; skewness falls from portfolio~1 to~5 for all countries but not monotonically for developed ones (Table~I, p.~45). Their pricing result is about volatility, not skewness: carry returns load negatively on innovations in global FX volatility, which price the five portfolios, while they ``do not find that the skewness of a portfolio is robustly related to average excess returns'' and aggregate-skewness betas are insignificant with low cross-sectional fit (sec.~V.C, pp.~28--29).

So negative skewness of carry is confirmed in this sample, including 2008, but it is not what explains the cross-section there. The Sharpe ratios are consistent with the means and standard deviations (7.23/9.81 = 0.74, 5.72/10.24 = 0.56). ``Measuring the standard error of the sample and finding the most attractive pairs'' is Gliner's (p.~209), without method.

\qffield{Pitfalls} Gliner's passage gives no evidence that the strategy outperforms standard carry trades. His Table~7.6 labels the carry ``3M Carry (LIBOR)'' and the volatility ``30D Volatility'', without saying that the volatility is realised.

\qffield{See also} \hyperref[meth:fx-valuation-signals]{FX valuation signals}; \hyperref[def:us-dollar-dxy-sdr]{the dollar}; \hyperref[meth:macro-position-sizing-risk]{gap risk}.

\qffield{Further reading} Carry crashes \citep{catalogue-0241}; carry and global volatility \citep{wide-menkhoff-carry}; \citep{catalogue-0240} on common risk factors in currency returns.

\end{qfmodel}
